\begin{scanpage}{0095}{78}
\[\left.\begin{aligned}E_\varphi=-H_\vartheta&=\frac{ia}2E_{0x}\frac{e^{ikR}}R\overline\Sigma_c^0(\vartheta,\gamma),\\E_\vartheta=H_\varphi&=\frac{ia}2H_{0x}\frac{e^{ikR}}R\Sigma_c^0(\vartheta,\gamma).\end{aligned}\right\}\tag{13.20}\]
Here
\[\overline\Sigma_c^0=G\sin\gamma,\quad\Sigma_c^0=-G\sin\vartheta,\tag{13.21}\]
and
\[\begin{aligned}G={}&2\sqrt{\frac2{\pi ka(\sin\gamma+\sin\vartheta)}}\times\\&\times\frac{\sin[\frac{kl}2(\cos\gamma-\cos\vartheta)]}{\cos\gamma-\cos\vartheta}e^{-ika(\sin\gamma+\sin\vartheta)+i\pi/4}.\end{aligned}\tag{13.22}\]
The index ``0'' on $\overline\Sigma_c^0$ and $\Sigma_c^0$ means that the field was calculated in the physical optics approach (based on the uniform part of the current), and the index ``c'' shows that this fringing field is created by a cylindrical surface. The effective scattering area, in accordance with (12.17), is determined for a cylindrical surface by the relationships
\[\left.\begin{aligned}\sigma_{c,E}^0&=\pi a^2|\overline\Sigma_c^0|^2=\pi a^2\sin^2\gamma|G|^2,\\\sigma_{c,H}^0&=\pi a^2|\Sigma_c^0|^2=\pi a^2\sin^2\vartheta|G|^2.\end{aligned}\right\}\tag{13.23}\]
In the direction of the mirror-reflected ray $(\vartheta=\gamma)$, we have
\[\sigma_{c,E}^0=\sigma_{c,H}^0=kal^2\sin\vartheta=\frac{2\pi a}\lambda l^2\sin\vartheta.\tag{13.24}\]
In the direction toward the source $(\vartheta=\pi-\gamma)$, the functions $\overline\Sigma_c^0$ and $\Sigma_c^0$ equal
\[\overline\Sigma_c^0=-\Sigma_c^0=\sqrt{\frac{\sin\vartheta}{\pi ka}}\frac{\sin(kl\cos\vartheta)}{\cos\vartheta}\cdot e^{-2ika\sin\vartheta+i\pi/4}.\tag{13.25}\]
These expressions are valid if $ka\sin\vartheta\gg1$. It is not difficult to see, by means of equations (13.02)--(13.05), that the fringing field equals zero if $\gamma=0$ and $\vartheta=\pi$. Thus in the case of radar (that
\end{scanpage}
